\chapter{Adaptation and Fine-Tuning}

\section*{학습 목표}

범용 VLA가 있다고 해서 곧바로 모든 연구실의 로봇을 잘 움직이는 것은 아니다. robot마다 camera pose, gripper, action frequency, joint limit, calibration, table height, object distribution, controller wrapper가 다르다. 따라서 adaptation과 fine-tuning은 선택 사항이 아니라 deployment process의 일부다. 이 장은 foundation VLA를 특정 robot/task/action space에 맞추는 engineering recipe를 정리한다.

\begin{enumerate}
  \item VLA adaptation이 단순 마지막 layer 학습이 아님을 이해한다.
  \item prompt, sensor, action head, adapter, LoRA, full fine-tuning의 차이를 구분한다.
  \item embodiment/action normalization과 catastrophic forgetting 문제를 정식화한다.
  \item offline evaluation에서 real-robot validation으로 넘어가는 protocol을 설계한다.
\end{enumerate}

\section{Why adaptation is unavoidable}

pretrained VLA policy를 다음과 같이 쓰자.
\begin{equation}
  \act_t
  \sim
  \pi_{\theta_0}
  (\cdot \mid \obs_{\leq t},\ell,\conf_t,m_0).
  \label{eq:ch15-pretrained-policy}
\end{equation}
여기서 $m_0$는 pretraining data의 implicit metadata이다. 실제 배치 대상 robot은 $m_{\star}$를 갖는다. adaptation은 다음 mismatch를 줄이는 과정이다.
\begin{equation}
  \Delta_m
  =
  d(m_0,m_{\star})
  =
  d_{sensor}
  +
  d_{action}
  +
  d_{embodiment}
  +
  d_{task}
  +
  d_{controller}.
  \label{eq:ch15-metadata-mismatch}
\end{equation}
이 mismatch가 작으면 prompt나 action normalization만으로 충분할 수 있다. mismatch가 크면 action head, vision encoder, language backbone, controller wrapper까지 조정해야 한다.

adaptation을 너무 약하게 하면 robot-specific behavior를 얻지 못한다. 너무 강하게 하면 pretrained representation을 망가뜨리거나 unsafe overfitting이 생긴다. 이 장의 핵심은 어디를 얼마나 바꿀지 결정하는 것이다.

\section{Adaptation targets}

VLA adaptation target은 여러 층에 있다.
\begin{table}[h]
  \centering
  \caption{VLA adaptation target}
  \begin{tabularx}{\linewidth}{p{0.22\linewidth}X X}
    \toprule
    Target & 무엇을 바꾸는가 & 적합한 상황 \\
    \midrule
    Prompt/template & instruction style, task wording & model capability는 충분하지만 language convention이 다를 때 \\
    Sensor remapping & camera crop, calibration, proprioception channel & observation schema가 다를 때 \\
    Action wrapper & normalization, clipping, frame convention & action space는 유사하지만 scale/rate가 다를 때 \\
    Action head & token/vector/chunk decoder & robot action type이 달라질 때 \\
    Adapter/LoRA & 일부 low-rank parameter update & data가 작고 pretrained representation을 보존해야 할 때 \\
    Full fine-tuning & backbone까지 update & domain gap이 크고 충분한 data/evaluation이 있을 때 \\
    Controller wrapper & safety, tracking, smoothing, rate conversion & model output과 robot controller가 직접 맞지 않을 때 \\
    \bottomrule
  \end{tabularx}
  \label{tab:ch15-adaptation-targets}
\end{table}

Open-source VLA 생태계에서 Octo, OpenVLA, OpenVLA-OFT, FAST, VLA Foundry 같은 흐름은 이 target들을 서로 다른 방식으로 다룬다 \cite{octo2024,openvla2024,openvlaoft2025,fast2025,vlafoundry2026}. 이 책에서는 이들을 특정 성능 수치의 증거가 아니라, adaptation design space를 보여주는 case study로 취급한다.

\section{Action head adaptation}

Chapter 13의 관점에서 가장 중요한 fine-tuning 대상은 action head이다. pretrained representation $h_t$가 주어졌을 때 action head는 다음 mapping을 만든다.
\begin{equation}
  \act_t
  =
  g_{\phi}(h_t,\conf_t).
  \label{eq:ch15-action-head}
\end{equation}
robot이 바뀌면 $g_{\phi}$의 output type이 달라질 수 있다. 예를 들어 source model이 token action을 내는데 target robot은 continuous delta pose를 요구할 수 있다.
\begin{equation}
  g_{\phi}^{source}: h_t \mapsto z_t,
  \qquad
  g_{\phi}^{target}: h_t \mapsto \Delta x_t.
  \label{eq:ch15-head-replacement}
\end{equation}
이 경우 full backbone보다 action head와 wrapper를 먼저 조정하는 것이 보수적이다. backbone은 semantic perception을 유지하고, target-specific motor interface만 바꾼다.

\section{Embodiment adapter}

multi-embodiment adaptation에서는 robot metadata를 policy에 명시적으로 넣는 것이 유리하다.
\begin{equation}
  e_{\star}
  =
  E_{\eta}(m_{\star}),
  \qquad
  \act_t
  \sim
  \pi_{\theta}
  (\cdot\mid \obs_{\leq t},\ell,\conf_t,e_{\star}).
  \label{eq:ch15-embodiment-adapter}
\end{equation}
여기서 $E_{\eta}$는 embodiment adapter이다. 이 adapter는 kinematics, action scale, camera calibration, gripper type 같은 정보를 representation에 넣는다.

adapter가 없으면 model은 robot-specific difference를 observation distribution에서 암묵적으로 추론해야 한다. 이는 data가 많을 때는 가능할 수 있지만, small fine-tuning에서는 불안정하다.

\section{LoRA, adapter tuning, and full fine-tuning}

parameter-efficient tuning의 한 예는 low-rank update이다.
\begin{equation}
  W'
  =
  W
  +
  \Delta W,
  \qquad
  \Delta W
  =
  BA,
  \quad
  \mathrm{rank}(\Delta W)\ll \mathrm{rank}(W).
  \label{eq:ch15-lora}
\end{equation}
이 방식은 pretrained representation을 크게 바꾸지 않고 target data에 적응할 수 있다. small dataset에서는 full fine-tuning보다 안정적일 수 있다.

하지만 parameter-efficient tuning이 항상 충분한 것은 아니다. camera modality가 크게 바뀌거나, action representation이 바뀌거나, target task가 pretraining distribution 밖이면 action head와 일부 encoder를 함께 조정해야 할 수 있다.

\begin{table}[h]
  \centering
  \caption{Fine-tuning 방식의 trade-off}
  \begin{tabularx}{\linewidth}{p{0.2\linewidth}X X}
    \toprule
    Method & 장점 & 위험 \\
    \midrule
    Frozen backbone + new head & 안정적이고 data-efficient & semantic-action mismatch가 남을 수 있음 \\
    Adapter / LoRA & forgetting을 줄이고 비용이 낮음 & 큰 domain gap에는 부족할 수 있음 \\
    Partial fine-tuning & 필요한 module만 바꿀 수 있음 & 어떤 layer를 열지 결정이 어려움 \\
    Full fine-tuning & target domain에 강하게 적응 & catastrophic forgetting, unsafe overfitting, 큰 검증 비용 \\
    \bottomrule
  \end{tabularx}
  \label{tab:ch15-ft-methods}
\end{table}

\section{Fine-tuning objective}

target dataset $\mathcal{D}_{target}$에 대해 behavior cloning loss를 쓰면
\begin{equation}
  \mathcal{L}_{target}
  =
  -
  \mathbb{E}_{(\obs,\ell,\act)\sim\mathcal{D}_{target}}
  \log
  \pi_{\theta}(\act\mid\obs,\ell).
  \label{eq:ch15-target-loss}
\end{equation}
그러나 target data가 작으면 overfitting과 forgetting이 생긴다. 따라서 source policy와의 regularization을 추가할 수 있다.
\begin{equation}
  \mathcal{L}
  =
  \mathcal{L}_{target}
  +
  \lambda
  D_{KL}
  (
  \pi_{\theta}(\cdot\mid\obs,\ell)
  \|
  \pi_{\theta_0}(\cdot\mid\obs,\ell)
  )
  +
  \beta
  \mathcal{L}_{safe}.
  \label{eq:ch15-regularized-ft}
\end{equation}
여기서 $\mathcal{L}_{safe}$는 unsafe action penalty, constraint violation penalty, 또는 safety classifier loss가 될 수 있다.

\section{Data mixture for adaptation}

adaptation data는 success-only demonstration으로 충분하지 않다.
\begin{equation}
  \mathcal{D}_{adapt}
  =
  \mathcal{D}_{demo}
  \cup
  \mathcal{D}_{correction}
  \cup
  \mathcal{D}_{failure}
  \cup
  \mathcal{D}_{intervention}.
  \label{eq:ch15-adaptation-data}
\end{equation}
small dataset에서는 data weighting이 중요하다.
\begin{equation}
  \mathcal{L}_{adapt}
  =
  \sum_{j}
  w_j
  \mathbb{E}_{r\sim\mathcal{D}_j}
  [
  \ell_{\theta}(r)
  ].
  \label{eq:ch15-adaptation-mixture}
\end{equation}
failure와 correction data는 수가 적어도 target deployment에서는 큰 영향을 줄 수 있다. 특히 action chunking이나 continuous action objective를 쓰는 경우, correction data는 recovery behavior를 학습하는 핵심이다.

\section{Action normalization and calibration}

fine-tuning에서 action normalization은 단순 preprocessing이 아니다. target robot의 physical scale을 정하는 calibration step이다.
\begin{equation}
  \act_t^{phys}
  =
  \sigma_{\star}
  \odot
  \act_t^{norm}
  +
  \mu_{\star},
  \qquad
  \act_t^{safe}
  =
  \mathrm{Clip}
  (\act_t^{phys},\mathcal{U}_{safe}).
  \label{eq:ch15-action-calibration}
\end{equation}
이 식에서 $\mu_{\star},\sigma_{\star}$와 $\mathcal{U}_{safe}$는 robot-specific이다. 잘못된 scale은 성공률 문제 이전에 safety issue가 된다.

calibration은 반드시 offline sanity check와 real robot low-speed test를 거쳐야 한다.
\begin{equation}
  \mathrm{Ready}
  =
  \mathbb{I}
  [
  \mathrm{SchemaCheck}
  \land
  \mathrm{ScaleCheck}
  \land
  \mathrm{SafetyCheck}
  \land
  \mathrm{LatencyCheck}
  ].
  \label{eq:ch15-readiness}
\end{equation}

\section{Validation protocol}

fine-tuning된 policy는 곧바로 full-speed real robot에 올리면 안 된다. 최소 세 단계가 필요하다.
\begin{enumerate}
  \item offline validation: held-out trajectory, action distribution, invalid action rate, latency를 본다.
  \item simulation or replay validation: sensor/action wrapper와 controller interface를 검증한다.
  \item guarded real-robot validation: low speed, restricted workspace, human intervention, emergency stop, logging을 사용한다.
\end{enumerate}

성공률 하나만으로는 충분하지 않다.
\begin{align}
  M_{deploy}
  =
  (&R_{success},
  R_{invalid},
  R_{saturation},
  T_{infer}, \notag\\
  &
  E_{tracking},
  R_{intervention},
  R_{recovery}).
  \label{eq:ch15-deployment-metrics}
\end{align}
이 지표들은 Chapter 16의 evaluation protocol로 이어진다.

\section{Case-study roles}

이 장에서 최신 오픈소스/프로젝트 사례는 다음 역할로 읽는다.
\begin{table}[h]
  \centering
  \caption{Adaptation ecosystem의 case-study 역할}
  \begin{tabularx}{\linewidth}{p{0.2\linewidth}X X}
    \toprule
    Case & 이 장에서의 역할 & 읽을 때의 주의점 \\
    \midrule
    Octo & generalist robot policy와 빠른 sensor/action adaptation의 사례 & dataset/action schema가 target robot과 맞는지 확인 \\
    OpenVLA & open-source VLA backbone과 robot action fine-tuning의 기준점 & token/action interface와 latency를 함께 봐야 함 \\
    OpenVLA-OFT & continuous action, chunking, practical fine-tuning recipe의 사례 & project/report claim과 independent validation 구분 \\
    FAST & efficient action tokenization과 compression의 사례 & tokenizer가 physical frequency structure를 보존하는지 확인 \\
    VLA Foundry & LLM/VLM/VLA training stack 통합의 engineering 사례 & framework capability와 model capability를 구분 \\
    \bottomrule
  \end{tabularx}
  \label{tab:ch15-case-study}
\end{table}

이름을 외우는 것보다 중요한 것은 각 사례가 어떤 adaptation boundary를 건드리는지 보는 것이다.

\begin{casebox}{Case study: OpenVLA-OFT as action-interface adaptation}
OpenVLA-OFT는 단순히 OpenVLA를 더 학습한 사례가 아니라 continuous action, action chunking, decoding speed, fine-tuning objective가 target robot deployment claim을 어떻게 바꾸는지 보여주는 사례다 \cite{openvlaoft2025}. Real VLA 관점의 질문은 성능 수치 하나가 아니라, target controller가 어떤 action을 받고 chunk interruption과 stale action을 어떻게 처리하는가이다.
\end{casebox}

\section{Deployment checklist}

\begin{enumerate}
  \item target robot의 action type, frame, frequency, controller를 명시했는가?
  \item source model의 action representation과 target action representation이 같은가?
  \item action normalization과 clipping이 robot-specific으로 검증되었는가?
  \item fine-tuning data에 failure, correction, intervention이 포함되어 있는가?
  \item LoRA/adapter/full fine-tuning 중 선택 이유가 data size와 domain gap에 맞는가?
  \item offline metric, replay/sim metric, guarded real metric을 분리했는가?
  \item deployment log가 다음 data engine loop로 들어가는가?
\end{enumerate}

\chapterwork
{Octo, OpenVLA, OpenVLA-OFT, FAST, VLA Foundry를 adaptation target과 evidence tier 관점에서 비교하라 \cite{octo2024,openvla2024,openvlaoft2025,fast2025,vlafoundry2026}.}
{pretrained VLA를 target robot에 맞출 때 sensor remapping보다 action normalization이 더 위험할 수 있는 이유는 무엇인가?}
{source VLA와 target robot이 다를 때 sensor adapter, action head, normalization, validation split, safety wrapper를 포함한 adaptation plan을 작성하라.}

\section{요약}

Adaptation은 VLA를 특정 robot system에 연결하는 과정이다. 이는 prompt engineering도 아니고 마지막 layer 학습만도 아니다. sensor remapping, action normalization, embodiment adapter, action head tuning, LoRA/full fine-tuning, data mixture, safety wrapper, validation protocol이 모두 포함된다. small data에서는 pretrained representation을 보존해야 하지만, target robot의 action/control boundary를 충분히 바꾸지 못하면 task는 실패한다. aggressive fine-tuning은 성능을 올릴 수 있지만 forgetting과 unsafe overfitting을 만든다.

다음 장에서는 fine-tuned VLA가 실제로 잘하는지 어떻게 주장할 수 있는지, benchmark와 evaluation protocol을 다룬다.
